\section{Evaluation}
\label{sec:eval}

\subsection{What is measured, and what is not}

Three things are measured: the size of the system, the size of what ships, and the wall-clock
cost of the operations the paper describes. All three are taken from the artefact on one
machine (Windows 11, Node 24, an x86-64 desktop CPU) with the scripts named in
\S\ref{sec:verification}.

We do \textbf{not} report a comparison against other readers, and the reason is not modesty.
A fair comparison would need a matched document corpus, a defined set of tasks, and a
measurement of correctness as well as speed --- and this system has no claim to be faster than
an engine written in C. What the numbers below establish is narrower and, we think, more
useful: that the writing operations are cheap enough to run on every save, and that
virtualisation makes open time independent of document length.

\subsection{Size}

\begin{table}[t]
  \centering
  \caption{What ships. The unpacked executable is the application plus the Electron runtime;
  the renderer bundle is the code this paper describes.}
  \label{tab:sizes}
  \small
  \begin{tabular}{@{}lr@{}}
    \toprule
    Artefact & Size \\
    \midrule
    Renderer bundle, JavaScript & 660.2\,kB (197.8\,kB gzip) \\
    Renderer bundle, CSS & 26.1\,kB (6.0\,kB gzip) \\
    UI resource container & 19.0\,kB \\
    Application source & 17{,}323 lines over 47 files \\
    Test and tooling source & 10{,}594 lines over 52 files \\
    \midrule
    Portable executable (single file) & 108.1\,MB \\
    Unpacked executable & 235.6\,MB \\
    \bottomrule
  \end{tabular}
\end{table}

Table~\ref{tab:sizes} separates two facts that are usually conflated. The code this paper
describes is 198\,kB compressed; the download is 108\,MB, and 99.8\,\% of it is the Electron
runtime. That ratio is a property of the packaging choice, not of the application, and it is
the project's largest engineering compromise: a native shell would ship a few megabytes and
cost months of platform work.

The test-to-application line ratio is 0.61:1. We report it because it is a design choice with
a cost --- a tenth of the codebase exists to check the other nine tenths --- and because a
figure that low is only defensible if the suites assert behaviour rather than implementation,
which \S\ref{sec:verification} argues they do.

\subsection{Cost of the writing operations}

\begin{table}[t]
  \centering
  \caption{Measured cost of the operations this paper describes, in the renderer, excluding
  file I/O. Median of repeated runs on fixtures of 24 and 200 pages; the raw samples are in
  \texttt{paper/\allowbreak metrics/\allowbreak benchmark.json}.}
  \label{tab:timings}
  \small
  \begin{tabular}{@{}lrr@{}}
    \toprule
    Operation & Document & Median \\
    \midrule
    Open to first painted page & 24 pages & 271\,ms \\
    Open to first painted page & 200 pages & 279\,ms \\
    Build full-text index & 24 pages & 4.9\,ms \\
    Build full-text index & 200 pages & 40.9\,ms \\
    Incremental annotation save & 24 pages & 1\,ms \\
    Whole-document rewrite & 24 pages & 2\,ms \\
    Merge two 24-page documents & 48 pages out & 3\,ms \\
    JavaScript heap after opening & 200 pages & 7\,MiB used \\
    \bottomrule
  \end{tabular}
\end{table}

Two rows carry the paper's practical argument.

\paragraph{Open time is independent of length.} 271\,ms for 24 pages and 279\,ms for 200 pages
--- a 3\,\% difference across an eight-fold increase in document size. This is what
virtualisation is for: the application rasterises the pages near the viewport and nothing else,
so the cost of opening is the cost of the first screen plus parsing the cross-reference, not
the cost of the document.

\paragraph{Writing is cheap enough to do on every save.} An incremental save of a 24-page
document is around a millisecond; a full rewrite of the same document is two; merging two
24-page documents into 48 is three. These are single-digit milliseconds for operations that
earlier sections described as delicate, which is the useful observation about the design: the
difficulty was never performance, and the invariants in \S\ref{sec:pdf-layer} cost nothing to
uphold. A save can therefore be conservative --- rewrite the page's annotation array, drop the
objects it replaces, write a fresh cross-reference --- rather than clever.

The index figure scales linearly: 4.9\,ms for 24 pages, 40.9\,ms for 200, i.e.\ about
0.2\,ms per page. That is why folder search rebuilds an index per file rather than storing one
(\S\ref{sec:folder-search}), and why the 7\,MiB heap figure for a 200-page document is
consistent with holding no more than the pages near the viewport.

\subsection{Threats to validity}

The measurements are from one machine and one rendering engine, and the electron and browser
versions are pinned only by \texttt{package.json} ranges. The timings are in-process and exclude
disk and IPC, so they understate what a user waits for on a save; the open-time figures include
navigation and first paint, so they overstate what a second open of the same file costs.

The assertion count is a measure of effort, not of coverage. There is no instrumentation-based
coverage figure in this paper because none was collected; a suite count can be raised by
writing easier tests, and the reader should treat 789 as evidence that many behaviours are
checked, not as a percentage of the behaviours that exist.
